%%%%%%%%%%%%%%%%%%%%%%%%%%%%%%%%%%%%%%%%%%%%%%%%%%%%%%%%%%%%%%%%%%%%%
%% This is a (brief) model paper using the achemso class
%% The document class accepts keyval options, which should include
%% the target journal and optionally the manuscript type. 
%%%%%%%%%%%%%%%%%%%%%%%%%%%%%%%%%%%%%%%%%%%%%%%%%%%%%%%%%%%%%%%%%%%%%
\documentclass[journal=jacsat,manuscript=article]{achemso}

%%%%%%%%%%%%%%%%%%%%%%%%%%%%%%%%%%%%%%%%%%%%%%%%%%%%%%%%%%%%%%%%%%%%%
%% Place any additional packages needed here.  Only include packages
%% which are essential, to avoid problems later. Do NOT use any
%% packages which require e-TeX (for example etoolbox): the e-TeX
%% extensions are not currently available on the ACS conversion
%% servers.%%%%%%%%%%%%%%%%%%%%%%%%%%%%%%%%%%%%%%%%%%%%%%%%%%%%%%%%%%%%%%%%%%%%%


 

%%%%%%%%%%%%%%%%%%%%%%%%%%%%%%%%%%%%%%%%%%%%%%%%%%%%%%%%%%%%%%%%%%%%%
\usepackage[version=3]{mhchem} % Formula subscripts using \ce{}

%%%%%%%%%%%%%%%%%%%%%%%%%%%%%%%%%%%%%%%%%%%%%%%%%%%%%%%%%%%%%%%%%%%%%
%% If issues arise when submitting your manuscript, you may want to
%% un-comment the next line.  This provides information on the
%% version of every file you have used.
%%%%%%%%%%%%%%%%%%%%%%%%%%%%%%%%%%%%%%%%%%%%%%%%%%%%%%%%%%%%%%%%%%%%%
%%\listfiles

%%%%%%%%%%%%%%%%%%%%%%%%%%%%%%%%%%%%%%%%%%%%%%%%%%%%%%%%%%%%%%%%%%%%%
%% Place any additional macros here.  Please use \newcommand* where
%% possible, and avoid layout-changing macros (which are not used
%% when typesetting).
%%%%%%%%%%%%%%%%%%%%%%%%%%%%%%%%%%%%%%%%%%%%%%%%%%%%%%%%%%%%%%%%%%%%%
\newcommand*\mycommand[1]{\texttt{\emph{#1}}}

\usepackage[version=3]{mhchem} % Formula subscripts using \ce{}
\usepackage{subcaption}
\usepackage{graphicx}
\usepackage[utf8]{inputenc}
\usepackage[english]{babel}
\usepackage{amsmath}
\usepackage{bm}
\usepackage{booktabs}
\usepackage{multirow}
\usepackage{pdflscape}
\usepackage{afterpage}
\usepackage{float}
\usepackage{xcolor}
\usepackage[symbol]{footmisc}
\renewcommand{\thefootnote}{\fnsymbol{footnote}}
\usepackage[T1, OT1]{fontenc}
\DeclareTextSymbolDefault{\DH}{T1}

\newcommand{\N}[0]{\mathcal{N}}
\newcommand{\Se}[0]{\mathcal{S}}
\newcommand{\M}[0]{\mathcal{M}}
\newcommand{\SM}[1]{\mathcal{S}_{M_{#1}}}


%%%%%%%%%%%%%%%%%%%%%%%%%%%%%%%%%%%%%%%%%%%%%%%%%%%%%%%%%%%%%%%%%%%%%
%% Meta-data block
%% ---------------
%% Each author should be given as a separate \author command.
%%
%% Corresponding authors should have an e-mail given after the author
%% name as an \email command. Phone and fax numbers can be given
%% using \phone and \fax, respectively; this information is optional.
%%
%% The affiliation of authors is given after the authors; each
%% \affiliation command applies to all preceding authors not already
%% assigned an affiliation.
%%
%% The affiliation takes an option argument for the short name.  This
%% will typically be something like "University of Somewhere".
%%
%% The \altaffiliation macro should be used for new address, etc.
%% On the other hand, \alsoaffiliation is used on a per author basis
%% when authors are associated with multiple institutions.
%%%%%%%%%%%%%%%%%%%%%%%%%%%%%%%%%%%%%%%%%%%%%%%%%%%%%%%%%%%%%%%%%%%%%

\author{Samuel Tovey}
\author{Marko Zecevic}
\author{Christian Holm}
\email{holm@icp.uni-stuttgart.de}
\affiliation[University of Stuttgart]
{Institute for Computational Physics, University of Stuttgart, Stuttgart, Germany}


%%%%%%%%%%%%%%%%%%%%%%%%%%%%%%%%%%%%%%%%%%%%%%%%%%%%%%%%%%%%%%%%%%%%%
%% The document title should be given as usual. Some journals require
%% a running title from the author: this should be supplied as an
%% optional argument to \title.
%%%%%%%%%%%%%%%%%%%%%%%%%%%%%%%%%%%%%%%%%%%%%%%%%%%%%%%%%%%%%%%%%%%%%
\title[Understanding Particle Correlation: Insights from Random Matrix Theory]
  {Understanding Particle Correlation: Insights from Random Matrix Theory}

%%%%%%%%%%%%%%%%%%%%%%%%%%%%%%%%%%%%%%%%%%%%%%%%%%%%%%%%%%%%%%%%%%%%%
%% Some journals require a list of abbreviations or keywords to be
%% supplied. These should be set up here, and will be printed after
%% the title and author information, if needed.
%%%%%%%%%%%%%%%%%%%%%%%%%%%%%%%%%%%%%%%%%%%%%%%%%%%%%%%%%%%%%%%%%%%%%
\abbreviations{MD,RMT,VN,GK}
\keywords{molecular dynamics, random matrix theory, transport coefficients, correlation matrix, von Neumann entropy}

%%%%%%%%%%%%%%%%%%%%%%%%%%%%%%%%%%%%%%%%%%%%%%%%%%%%%%%%%%%%%%%%%%%%%
%% The manuscript does not need to include \maketitle, which is
%% executed automatically.
%%%%%%%%%%%%%%%%%%%%%%%%%%%%%%%%%%%%%%%%%%%%%%%%%%%%%%%%%%%%%%%%%%%%%
\begin{document}

%%%%%%%%%%%%%%%%%%%%%%%%%%%%%%%%%%%%%%%%%%%%%%%%%%%%%%%%%%%%%%%%%%%%%
%% The "tocentry" environment can be used to create an entry for the
%% graphical table of contents. It is given here as some journals
%% require that it is printed as part of the abstract page. It will
%% be automatically moved as appropriate.
%%%%%%%%%%%%%%%%%%%%%%%%%%%%%%%%%%%%%%%%%%%%%%%%%%%%%%%%%%%%%%%%%%%%%
\begin{tocentry}

\end{tocentry}

%%%%%%%%%%%%%%%%%%%%%%%%%%%%%%%%%%%%%%%%%%%%%%%%%%%%%%%%%%%%%%%%%%%%%
%% The abstract environment will automatically gobble the contents
%% if an abstract is not used by the target journal.
%%%%%%%%%%%%%%%%%%%%%%%%%%%%%%%%%%%%%%%%%%%%%%%%%%%%%%%%%%%%%%%%%%%%%

\begin{abstract}
 Particle correlations underpin a wide range of calculations in molecular dynamics and scattering experiments, most visibly through the Green--Kubo relations that connect microscopic fluctuations to macroscopic transport coefficients.
While correlation functions are routinely computed, the particle-wise matrix from which they are assembled has received comparatively little attention as an object of study in its own right.
In this work we examine this matrix directly, drawing on random matrix theory to separate physical structure from finite-sampling noise, and on spectral information measures to track its behaviour across simulation conditions.
Using theoretical arguments together with molecular dynamics simulations of a Lennard-Jones fluid, bonded systems, water, and molten salts, we show that the spectrum and the von Neumann entropy of the particle-wise velocity correlation matrix resolve bonding topology and first-order phase transitions, and in some cases do so more sharply than standard thermodynamic observables.

\end{abstract}
\newpage
%%%%%%%%%%%%%%%%%%%%%%%%%%%%%%%%%%%%%%%%%%%%%%%%%%%%%%%%%%%%%%%%%%%%%
%% Start the main part of the manuscript here.
%%%%%%%%%%%%%%%%%%%%%%%%%%%%%%%%%%%%%%%%%%%%%%%%%%%%%%%%%%%%%%%%%%%%%
    % Begin the document content
    \begin{table}[H]
\begin{tabular}{@{}ll@{}}
\toprule
Symbol                               & \multicolumn{1}{c}{Description}                                                                                                                                                                               \\ \midrule
                                     &                                                                                                                                                                                                               \\
$\mathcal{U}$                        & A finite system of atoms.                                                                                                                                                                                     \\
                                     &                                                                                                                                                                                                               \\
$\mathcal{S}$                        & The set of all unique atomic species in $\mathcal{U}$                                                                                                                                                         \\
                                     &                                                                                                                                                                                                               \\
$\mathcal{N}$                        & The indexing set of all atoms in $\mathcal{U}$                                                                                                                                                                \\
                                     &                                                                                                                                                                                                               \\
$\mathcal{N}_{\alpha}$               & \begin{tabular}[c]{@{}l@{}}The indices of the atoms of species $\alpha \in \mathcal{S}$ such that $|\mathcal{N}_{\alpha}|$ \\ is the total number of atoms of species $\alpha$ in $\mathcal{U}$.\end{tabular} \\
                                     &                                                                                                                                                                                                               \\
$\mathcal{M}$                        & The indexing set of all molecules in $\mathcal{U}$                                                                                                                                                            \\
                                     &                                                                                                                                                                                                               \\
$\mathcal{B}$                        & The set of all molecular species in $\mathcal{U}$                                                                                                                                                             \\
                                     &                                                                                                                                                                                                               \\
$\mathcal{M}_{\alpha}$               & The indices of the molecules of species $\alpha \in \mathcal{B}$                                                                                                                                              \\
                                     &                                                                                                                                                                                                               \\
$\mathcal{S}_{\mathcal{M}_{\alpha}}$ & \begin{tabular}[c]{@{}l@{}}Set of all atomic species from which molecule $\alpha \in \mathcal{B}$ is \\ constructed. Therefore, $\mathcal{S}_{\mathcal{M}_{\alpha}} \subset \mathcal{S}$\end{tabular}         \\
                                     &                                                                                                                                                                                                               \\
$\mathcal{M}^{i}_{\alpha}$           & \begin{tabular}[c]{@{}l@{}}Indices of the atoms belonging to the $i^{\text{th}}$ molecule of \\ species $\alpha \in \mathcal{B}$.\end{tabular}                                                                \\
                                     &                                                                                                                                                                                                               \\
$\mathcal{M}^{i}_{\alpha \beta}$     & \begin{tabular}[c]{@{}l@{}}Indices of the atoms of species $\beta \in \mathcal{S}_{\mathcal{M}_{\alpha}}$ belonging to the $i^{\text{th}}$  \\ molecule of  species $\alpha \in \mathcal{B}$\end{tabular}     \\
                                     &                                                                                                                                                                                                               \\
$f_{\alpha}$                         & $f_{\alpha} = \int\limits_{0}^{\infty} dt \sum\limits_{i \in \mathcal{N}_{\alpha}} \langle \mathbf{v}_{i}(t) \cdot \mathbf{v}_{i}(0) \rangle$                                                                 \\
                                     &                                                                                                                                                                                                               \\
$f_{\alpha \beta}$                   & $f_{\alpha \beta} = \int\limits_{0}^{\infty} dt \sum\limits_{i \in \mathcal{N}_{\alpha}}\sum\limits_{ j \in \mathcal{N}_{\beta}} \langle \mathbf{v}_{i}(t) \cdot \mathbf{v}_{j}(0) \rangle$                   \\
                                     &                                                                                                                                                                                                               \\
$g_{\alpha \beta}$                   & $g_{\alpha \beta} = f_{\alpha\beta} - \delta_{\alpha\beta}\cdot f_{\alpha}$                                                                                                                                   \\
                                     &                                                                                                                                                                                                               \\ \bottomrule
\end{tabular}
\caption{Mathematical definitions used throughout the paper.}
\label{tab:definitions}
\end{table}
    \newpage
    \section{Introduction}
\label{sec:introduction}

Time-correlation functions occupy a central position in statistical mechanics.
Through the Green--Kubo relations~\cite{green54a,kubo57a,kubo66a}, equilibrium fluctuations of a microscopic observable determine macroscopic transport coefficients; through the related Onsager reciprocal relations~\cite{onsager31a,onsager31b}, those coefficients exhibit the symmetries required by microscopic reversibility; and through van Hove correlation functions and their reciprocal-space transforms, they are tied to scattering experiments on dense fluids and soft matter~\cite{hansen90a}.
For essentially any observable derivable from the dynamics of a many-body system --- diffusion, ionic conductivity, shear viscosity, thermal conductivity, density fluctuations --- a correlation function can be written down, measured in simulation, and compared against experiment.

In practice, however, the information contained in these correlation functions is almost always compressed before it is used.
A viscosity is reported as a scalar; an ionic conductivity is reported as a scalar, sometimes decomposed into a Nernst--Einstein baseline and a distinct-correlation correction tailored to that specific coefficient~\cite{kashyap11a,fong20a,fong21a}; a dynamic structure factor is reported as a two-point function of wavevector and frequency.
Each of these reductions discards the particle-wise matrix structure from which it was assembled, and the self-versus-distinct decompositions that do exist are written down property-by-property rather than as a single framework that any transport coefficient of the relevant form can be dropped into.
Working backwards from such a coefficient to its microscopic definition recovers a sum over pairs of particles of an integrated velocity (or current) cross-correlation, and that object --- an $|\N| \times |\N|$ real symmetric matrix over all atoms in the system --- is seldom examined on its own terms.

This omission is not for lack of tools.
Correlation matrices constructed from high-dimensional data are studied extensively in other settings, where random matrix theory (RMT) provides a natural language for distinguishing physical structure from finite-sampling noise~\cite{TODO:plerou1999,TODO:laloux1999,TODO:bouchaud2009}.
In the molecular-dynamics literature, the closest relatives are \emph{configurational} in nature: Schlitter's formula~\cite{TODO:schlitter1993} and its refinements~\cite{TODO:andricioaei2001} compute an entropy from the spectrum of the mass-weighted positional covariance matrix, and essential-dynamics~\cite{TODO:amadei1993} / quasi-harmonic analyses diagonalise the same object to extract collective modes.
The matrix assembled from integrated velocity cross-correlations --- the object that actually appears in the Green--Kubo decomposition of transport coefficients --- has, to our knowledge, not been treated with the same spectral machinery.

In this work we do so, and in doing so make a contribution in two parts.
First, we give a species-generic Green--Kubo decomposition that expresses any transport coefficient of the form $\lambda \propto \int_{0}^{\infty} dt \, \langle A(t) \cdot A(0) \rangle$, with $A = \sum_{i} \eta_{\alpha(i)}\, a_{i}$ built from a single-particle primitive $a_{i}$ and a species-specific scalar $\eta_{\alpha}$, as a coarse linear functional of a single underlying object: the atom-wise time-integrated correlation matrix $C$.
This places self-diffusion, ionic conductivity, electrophoretic mobility, and related dissipative quantities on a common footing, rather than re-deriving a bespoke self / distinct decomposition for each one.
Second, we study $C$ directly.
We use its eigenvalue spectrum to identify physical correlation structure; we use its von Neumann entropy as a scalar summary of how correlation is distributed across the system; and we develop a Marchenko--Pastur-type null against which the bulk of the spectrum can be compared, so that physical signal is separable from finite-sample noise.
Using Lennard-Jones fluids, a controlled bond-addition protocol, water across its condensed phases, and a molten salt, we show that eigenvalue outliers of $C$ are diagnostic of bonding topology and that the von Neumann entropy of $C$ resolves first-order phase transitions at least as sharply as conventional thermodynamic observables, and in some cases more sharply.

The remainder of the paper is organised as follows.
Section~\ref{sec:theory} reviews the Onsager and Green--Kubo framework in a form suited to species-wise decomposition, introduces the matrix $C$ and the spectral quantities we use to characterise it, and states the RMT null.
Section~\ref{sec:methods} describes the simulation protocols, the construction of $C$ from a trajectory, and the noise-correction procedure.
Section~\ref{sec:experiments} presents the four experimental studies described above.

    \section{Theory}
\label{sec:theory}

\subsection{Onsager and Green--Kubo Relations}
\label{subsec:onsager-relations}
Consider a system close to thermodynamic equilibrium in which a small set of thermodynamic forces $\{X_{j}\}$ drives conjugate fluxes $\{J_{i}\}$.
Onsager's linear-response ansatz~\cite{onsager31a,onsager31b} relates the averaged fluxes and forces through a transport-coefficient matrix $\mathbf{L}$,
\begin{align}
    \langle J_{i} \rangle = \sum\limits_{j} L_{ij}\, X_{j},
    \label{eqn:onsager-linear-response}
\end{align}
with the reciprocal relations $L_{ij} = L_{ji}$ following from microscopic reversibility~\cite{callen51a}.
The fluctuation--dissipation theorem, in the form written down by Green~\cite{green54a} and Kubo~\cite{kubo57a,kubo66a}, gives each $L_{ij}$ as the time integral of an equilibrium flux--flux correlation function,
\begin{align}
    L_{ij} = \frac{1}{k_{B}T V} \int\limits_{0}^{\infty} dt \, \langle J_{i}(t) \cdot J_{j}(0) \rangle,
    \label{eqn:green-kubo-general}
\end{align}
where $V$ is the system volume and $T$ its temperature.
Equation~\ref{eqn:green-kubo-general} is the starting point for computing transport coefficients from molecular dynamics trajectories and is the relation we decompose into particle-wise contributions in the remainder of this section.

\subsection{Microscopic Decomposition of Transport Coefficients}
\label{subsec:microscopic-origins-of-transport-properties}
The Green--Kubo relations of Section~\ref{subsec:onsager-relations} can be used to compute a wide range of observables, from particle-specific quantities such as self-diffusion to system-wide quantities such as ionic conductivity.
Nernst--Einstein-type decompositions of such coefficients into a self contribution and a correction from distinct pair correlations are well established for individual transport properties, most prominently for ionic conductivity in electrolytes and molten salts~\cite{kashyap11a,fong20a,fong21a}.
Those derivations are, however, written down coefficient by coefficient; to our knowledge they have not been cast in a species-generic form that treats on a single footing any dissipative quantity derivable from a single-particle primitive.
We do so here.
The resulting formalism makes the \emph{same} distinct-correlation correction $\Delta^{\lambda}$ --- and, as we will show in Section~\ref{subsec:atom-wise-correlation-matrix}, the \emph{same} underlying matrix $C$ --- the common object behind self-diffusion, ionic conductivity, electrophoretic mobility, and the other transport coefficients built from a particle-wise observable with a species-specific scalar.
The notation used throughout this section is collected in Table~\ref{tab:definitions}.

Consider a dissipative quantity $\lambda$ related by a Green--Kubo expression to a system-wide observable $A$,
\begin{align}
    \lambda = \frac{\beta}{3} \int\limits_{0}^{\infty} dt \, \langle A(t)\cdot A(0) \rangle,
    \label{eqn:genericsystem}
\end{align}
where $\beta = 1/k_{B}T$.
As $A(t)$ and $A(0)$ refer to the same observable evaluated at different times, we call $\langle A(t)\cdot A(0)\rangle$ an autocorrelation function.
Suppose $A(t)$ decomposes into a sum of particle-wise contributions,
\begin{align}
    A(t) = \sum\limits_{i \in \N} a_{i}(t),
    \label{eqn:localproperties}
\end{align}
where $a_{i}(t)$ is a property associated with a single particle (e.g.\ velocity, position, or atomic current) in a system of $|\N|$ particles.
Equation~\ref{eqn:localproperties} may be further separated into same-species contributions,
\begin{align}
    A(t) = \sum\limits_{\alpha \in \Se} \sum\limits_{j \in \N_{\alpha}} a_{j}(t),
\end{align}
where $\alpha$ indexes a single species in the system.
Substituting this expansion into Equation~\ref{eqn:genericsystem} yields
\begin{align}
    \lambda = \frac{\beta}{3} \Bigg(
    &\int\limits_{0}^{\infty} dt \sum\limits_{\alpha \in \Se} \sum\limits_{i \in \N_{\alpha}} \langle a_{i}(t)\cdot a_{i}(0) \rangle + \nonumber \\
    &\int\limits_{0}^{\infty} dt \sum\limits_{\alpha \in \Se} \sum\limits_{\substack{i,j \in \N_{\alpha} \\ i \neq j}} \langle a_{i}(t)\cdot a_{j}(0) \rangle + \nonumber \\
    &\int\limits_{0}^{\infty} dt \sum\limits_{\substack{\alpha, \beta \in \Se \\ \alpha \neq \beta}} \sum\limits_{i \in \N_{\alpha}} \sum\limits_{j \in \N_{\beta}} \langle a_{i}(t)\cdot a_{j}(0) \rangle
    \Bigg),
    \label{eqn:expansion}
\end{align}
where the first term is the self contribution and the remaining two are the so-called distinct contributions, separated into same-species and cross-species parts.
To write the subsequent expressions compactly, we define
\begin{align}
    f_{\alpha}         &= \int\limits_{0}^{\infty} dt \sum\limits_{i \in \N_{\alpha}} \langle a_{i}(t)\cdot a_{i}(0) \rangle, \nonumber \\
    f_{\alpha \beta}   &= \int\limits_{0}^{\infty} dt \sum\limits_{i \in \N_{\alpha}} \sum\limits_{j \in \N_{\beta}} \langle a_{i}(t)\cdot a_{j}(0) \rangle, \nonumber \\
    g_{\alpha \beta}   &= f_{\alpha \beta} - \delta_{\alpha \beta}\, f_{\alpha},
    \label{eqn:fgdefs}
\end{align}
consistent with the definitions collected in Table~\ref{tab:definitions}.
Substituting Equation~\ref{eqn:fgdefs} into Equation~\ref{eqn:expansion} gives
\begin{align}
    \lambda = \frac{\beta}{3} \Bigg( \sum\limits_{\alpha \in \Se} f_{\alpha}
                                   + \sum\limits_{\alpha \in \Se} g_{\alpha \alpha}
                                   + \sum\limits_{\substack{\alpha, \beta \in \Se \\ \alpha \neq \beta}} g_{\alpha \beta} \Bigg),
    \label{eqn:simplification}
\end{align}
in which the self term and the same-species and cross-species distinct terms appear separately.
The effect of the distinct contributions can be absorbed into a dimensionless correction,
\begin{align}
    \lambda &= \lambda^{\text{NE}} \cdot (1 + \Delta^{\lambda}),
    \label{eqn:transport-equation}
\end{align}
where
\begin{align}
    \Delta^{\lambda} &= \frac{\sum\limits_{\beta \in \Se} \sum\limits_{\gamma \in \Se} \eta_{\beta}\, \eta_{\gamma}\, g_{\beta \gamma}}
                               {\sum\limits_{\alpha \in \Se} \eta_{\alpha}^{2}\, f_{\alpha}}.
    \label{eqn:full-correction}
\end{align}
Here we assume that the microscopic quantity $a_{i}$ can be separated into a dynamical part and a species-specific scalar $\eta_{\alpha}$, as is the case for charge or mass.
When $\Delta^{\lambda} = 0$ the expression reduces to the corresponding Nernst--Einstein-type estimate $\lambda^{\text{NE}}$, which ignores distinct correlations; $\Delta^{\lambda}$ therefore measures the deviation of the full transport coefficient from its uncorrelated single-particle baseline.
In some cases, while the transport property of interest is collective, it remains species-resolved --- electrophoretic mobility being one example.
Equations~\ref{eqn:transport-equation} and~\ref{eqn:full-correction} then specialise to
\begin{align}
    \lambda_{\alpha} &= \lambda_{\alpha}^{\text{NE}} \cdot (1 + \Delta_{\alpha}^{\lambda}),
\end{align}
with
\begin{align}
    \Delta_{\alpha}^{\lambda} &= \frac{\sum\limits_{\beta \in \Se} \eta_{\beta}\, g_{\alpha \beta}}{\eta_{\alpha}\, f_{\alpha}}.
\end{align}
This decomposition exposes the microscopic origin of transport coefficients and motivates treating the underlying correlation structure as an object of direct interest.

We note in closing a scope caveat.
The derivation above presumes that $A$ admits a single-particle expansion of the form of Equation~\ref{eqn:localproperties}, which is immediate for self-diffusion, ionic conductivity, and electrophoretic mobility.
For transport coefficients whose flux carries an explicit pair-interaction contribution --- shear viscosity and thermal conductivity being the canonical examples --- an Irving--Kirkwood-type atom-wise partition of the interaction term~\cite{TODO:irving-kirkwood} is required to bring the observable into the required form, after which the same decomposition applies.
We do not develop that partition here and restrict the experiments in Section~\ref{sec:experiments} to transport coefficients that admit the single-particle expansion directly.

\subsection{The Atom-Wise Correlation Matrix}
\label{subsec:atom-wise-correlation-matrix}
The quantities $f_{\alpha}$, $f_{\alpha \beta}$, and $g_{\alpha \beta}$ introduced in Equation~\ref{eqn:fgdefs} are scalar sums; the information in the summands is richer.
We therefore promote the particle indices to matrix indices and define the \emph{atom-wise correlation matrix}
\begin{align}
    C_{ij} = \int\limits_{0}^{\infty} dt \, \langle a_{i}(t) \cdot a_{j}(0) \rangle, \qquad i, j \in \N.
    \label{eqn:Cij}
\end{align}
By time-reversal symmetry of equilibrium correlation functions, $C$ is real and symmetric, with diagonal entries set by single-particle autocorrelations and off-diagonals carrying the distinct correlations that drive $\Delta^{\lambda}$.

Ordering the atom indices so that atoms of a given species appear contiguously partitions $C$ into species blocks: the $|\N_{\alpha}| \times |\N_{\beta}|$ block $C^{(\alpha \beta)}$ collects every correlation between an atom of species $\alpha$ and an atom of species $\beta$.
In this representation, the scalars of Equation~\ref{eqn:fgdefs} are block sums --- $f_{\alpha}$ is the trace of $C^{(\alpha \alpha)}$, $f_{\alpha \beta}$ is the sum of all entries of $C^{(\alpha \beta)}$, and $g_{\alpha \beta}$ is the off-trace sum --- so the decomposition in Equations~\ref{eqn:expansion}--\ref{eqn:simplification} is a statement about how the macroscopic transport coefficient $\lambda$ depends on coarse linear functionals of $C$.
Moving from those linear functionals to the spectrum of $C$ is the step that takes us beyond what a transport coefficient reports.

$C$ is itself a finite-sample estimator.
For a trajectory of length $T_{\text{traj}}$, Equation~\ref{eqn:Cij} is evaluated as a time-origin average, so the entries of $C$ carry sampling noise whose magnitude scales with the effective number of independent time samples $T_{\text{eff}} = T_{\text{traj}} / \tau_{\text{int}}$, where $\tau_{\text{int}}$ is the integral relaxation time of the underlying velocity autocorrelation.
Across an ensemble of independent trajectories, each realisation of $C$ is drawn from a distribution determined by the underlying dynamics, and $C$ is natural to regard as a random matrix.
In the limiting case of non-interacting particles, the $|\N|$ time series $\{a_{i}(t)\}$ are independent, so the off-diagonal entries of $C$ are inner products of independent signals and have zero mean; the resulting spectrum converges to a Marchenko--Pastur distribution~\cite{TODO:marchenko1967} set by $|\N|$ and $T_{\text{eff}}$.
This uncorrelated limit provides the null against which we measure deviation: physical signal --- bonded pairs, solvation shells, collective modes associated with phase behaviour --- appears as structure inconsistent with the noise-only reference, most directly as eigenvalue outliers above the bulk edge.

\subsection{Spectral Observables}
\label{subsec:interpreting-atomic-correlation-matrices}
We characterise $C$ through two families of observables: its eigenvalue spectrum and scalar summaries of that spectrum.
Because the integrated autocorrelation $\int_{0}^{\infty} dt\, \langle X(t)\cdot X(0) \rangle$ of any linear combination $X = \sum_{i} x_{i}\, a_{i}$ is non-negative in equilibrium~\cite{hansen90a}, $C$ is positive semi-definite and admits an orthonormal eigendecomposition $C = \sum_{k} \mu_{k}\, \phi_{k}\, \phi_{k}^{\top}$ with $\mu_{k} \ge 0$.

\subsubsection{Spectrum}
The eigenvalue distribution of $C$ separates naturally into a bulk and a set of outliers.
The bulk remains close to the Marchenko--Pastur null of Section~\ref{subsec:atom-wise-correlation-matrix} and encodes the sampling noise present even in an uncorrelated system.
Outliers above the bulk edge correspond to collective modes whose eigenvectors $\phi_{k}$ identify which atoms participate in a given correlated motion.
The participation ratio
\begin{align}
    \mathrm{PR}_{k} = \frac{1}{\sum_{i} \phi_{ik}^{4}}
    \label{eqn:participation-ratio}
\end{align}
counts, in a normalised sense, how many atoms are active in mode $k$: bonded pairs yield $\mathrm{PR} \approx 2$, small clusters yield $\mathrm{PR}$ of order the cluster size, and system-spanning collective modes yield $\mathrm{PR}$ of order $|\N|$.
This gives a direct route from matrix structure to physical interpretation: a large eigenvalue with small $\mathrm{PR}$ flags a local correlation such as a bond, whereas a large eigenvalue with large $\mathrm{PR}$ flags a long-range collective mode of the kind expected near a phase transition.

\subsubsection{Entropy}
The \emph{von Neumann entropy} of $C$ is defined as
\begin{align}
    S(C) = -\sum\limits_{k} p_{k} \ln p_{k}, \qquad p_{k} = \frac{\mu_{k}}{\sum_{l} \mu_{l}},
    \label{eqn:vnentropy}
\end{align}
where $\{\mu_{k}\}$ are the eigenvalues of $C$.
$S(C)$ is a scalar summary of how broadly correlation is distributed across the eigenmodes of $C$: it is maximal when eigenvalues are uniform (fully delocalised correlation) and minimal when a single eigenvalue dominates (one collective mode).
$S(C)$ can also be read as an entropy of the trace-normalised operator $C / \mathrm{Tr}(C)$, which makes explicit the formal parallel with the von Neumann entropy of a density matrix in quantum statistics and with the spectral configurational entropy of Schlitter~\cite{TODO:schlitter1993} applied to positional covariance.
We use $S(C)$ as the primary spectral observable throughout the experiments that follow.

    \section{Methods}
\label{sec:methods}

\subsection{Construction of the Correlation Matrix}
\label{subsec:matrix-computation}
Given a molecular dynamics trajectory of length $T_{\text{traj}}$ from which we sample $M$ velocity snapshots at uniform intervals $\delta t$, we compute the finite-sample estimator of Equation~\ref{eqn:Cij} as
\begin{align}
    \hat{C}_{ij} = \delta t \sum\limits_{\tau = 0}^{\tau_{\max}/\delta t} \frac{1}{M - \tau} \sum\limits_{m = 0}^{M - \tau - 1} \mathbf{v}_{i}(t_{m} + \tau\, \delta t) \cdot \mathbf{v}_{j}(t_{m}),
    \label{eqn:Chat}
\end{align}
where the outer sum is a Riemann approximation to the time integral, truncated at a lag $\tau_{\max}$ that we choose to be several times the integral relaxation time of the velocity autocorrelation (see below), and the inner sum averages over time origins $t_{m}$.

Evaluating Equation~\ref{eqn:Chat} by direct summation costs $\mathcal{O}(|\N|^{2} M^{2})$ and is intractable at the system sizes we use.
Instead, for each pair $(i,j)$ we compute the cross-correlation in frequency space via the Wiener--Khinchin theorem, using a single pair of FFTs per Cartesian component and summing the three component contributions, reducing the cost to $\mathcal{O}(|\N|^{2} M \log M)$.
The implementation is memory-bound at large $|\N|$, so we process blocks of pair indices in sequence and reuse the forward transforms of each atom's velocity time series.

Two practical points deserve attention.
First, when $a_{i}$ is a vector such as velocity, we compute $\hat{C}_{ij}$ per Cartesian component and average; this reduces variance by a factor of three relative to using a scalar projection and preserves rotational invariance.
Second, the choice of $\tau_{\max}$ is not innocuous.
Too short a cutoff under-integrates long-lived correlations and biases $\hat{C}$ downward; too long a cutoff accumulates noise from the flat tail of the autocorrelation function and inflates the variance of every off-diagonal entry.
We estimate the integral relaxation time
\begin{align}
    \tau_{\text{int}} = \int\limits_{0}^{\infty} dt \, \frac{\langle \mathbf{v}_{i}(t) \cdot \mathbf{v}_{i}(0) \rangle}{\langle \mathbf{v}_{i}(0) \cdot \mathbf{v}_{i}(0) \rangle}
    \label{eqn:tau-int}
\end{align}
directly from the trajectory, averaging over atoms, and set $\tau_{\max} = 10\, \tau_{\text{int}}$ by default.
The same estimate provides the effective number of independent time samples $T_{\text{eff}} = T_{\text{traj}} / \tau_{\text{int}}$ used in the noise model of Section~\ref{subsec:de-noising-correlation-matrices}.

\subsection{De-noising via the Marchenko--Pastur Null}
\label{subsec:de-noising-correlation-matrices}
For a system of $|\N|$ non-interacting particles sampled at $T_{\text{eff}}$ independent times, the bulk of the eigenvalue distribution of the trace-normalised correlation matrix $\hat{C} / \mathrm{Tr}(\hat{C})$ follows the Marchenko--Pastur density~\cite{TODO:marchenko1967},
\begin{align}
    \rho_{\text{MP}}(\mu) = \frac{1}{2 \pi q\, \sigma^{2}\, \mu} \sqrt{(\mu_{+} - \mu)(\mu - \mu_{-})}, \qquad \mu \in [\mu_{-}, \mu_{+}],
    \label{eqn:mp-density}
\end{align}
with aspect ratio $q = |\N| / T_{\text{eff}}$, edges $\mu_{\pm} = \sigma^{2}\, (1 \pm \sqrt{q})^{2}$, and noise scale $\sigma^{2}$ fitted from the bulk.
We refer to $\rho_{\text{MP}}$ as the noise-only null and use it in two ways.

First, we test how close the empirical spectrum $\rho_{\text{emp}}(\mu)$ is to the null using the Kolmogorov--Smirnov distance $D_{\mathrm{KS}}(\rho_{\text{emp}}, \rho_{\text{MP}})$ and a comparison between the empirical largest eigenvalue $\mu_{\max}$ and the MP edge $\mu_{+}$.
For a system with genuine correlation structure, $D_{\mathrm{KS}}$ remains non-trivial even in the $T_{\text{eff}} \to \infty$ limit and $\mu_{\max}$ sits above $\mu_{+}$; for an uncorrelated reference, both quantities approach zero.

Second, we use the null to separate signal from noise in the entropy observable.
Eigenvalues above a threshold $\mu_{\text{cut}} = \mu_{+} (1 + \kappa / \sqrt{T_{\text{eff}}})$, with $\kappa$ a small safety factor set to $\kappa = 3$ throughout, are classified as outliers; the remainder of the spectrum is taken to be consistent with the null.
We report two entropy quantities: the raw $S(\hat{C})$ of Equation~\ref{eqn:vnentropy} and the de-noised version
\begin{align}
    S_{\text{signal}}(\hat{C}) = S(\hat{C}) - S_{\text{bulk}}(\hat{C}),
    \label{eqn:signal-entropy}
\end{align}
where $S_{\text{bulk}}$ is the von Neumann entropy computed from only those eigenvalues below $\mu_{\text{cut}}$.
$S_{\text{signal}}$ is the component of the spectrum that cannot be explained by the noise-only null and is the quantity most sensitive to the physical features we want to expose.

We validate this machinery on synthetic trajectories with prescribed correlation structure before applying it to molecular dynamics output; the validation is reported in Experiment~\ref{subsec:convergence-of-the-noise-term}.

\subsection{Molecular Dynamics Simulations}
\label{subsec:molecular-dynamics-simulations}
Simulations fall into three families, one per experimental system, and use the classical molecular dynamics code LAMMPS~\cite{TODO:lammps} unless otherwise noted.
All production runs use Langevin dynamics in the constant-$NVT$ ensemble unless $NPT$ is explicitly required, with a damping time of $0.1$~ps.
For $NPT$ runs we use a Nos\'e--Hoover barostat with coupling constants of $0.1$~ps (thermostat) and $1.0$~ps (barostat).
Each run is initialised from a randomly seeded velocity distribution and separately reported in output metadata, and every reported observable is averaged over a minimum of three and up to five independent seeds.

\paragraph{Lennard-Jones fluids (Experiments~\ref{subsec:convergence-of-the-noise-term}, \ref{subsec:bond-addition}, and~\ref{subsec:lennard-jones-phases}).}
The Lennard-Jones potential is taken with argon parameters $\sigma = 3.4$~\AA{}, $\varepsilon / k_{B} = 120$~K, and particle mass $40$~amu, cut and shifted at $r_{c} = 2.5\,\sigma$.
System size is $|\N| = 512$ for the phase sweep and the bond-addition experiment, and we sweep $|\N| \in \{64, 256, 1024, 4096\}$ in the noise-convergence study.
The integration timestep is $\delta t = 2$~fs.
For the phase sweep we equilibrate for $10^{5}$ steps in $NPT$ at $P = 1$~atm and run for $5 \times 10^{5}$ steps of production; for the bond-addition experiment we equilibrate an $NVT$ liquid at $T = 80$~K and $\rho = 0.8\,\sigma^{-3}$ before inserting bonds and running for $5 \times 10^{5}$ production steps; for the non-interacting noise study we disable the pair potential entirely and run at low density and high temperature.
Harmonic bonds for the bond-addition experiment are added between random eligible pairs with stiffness $k_{b} = 100\,\varepsilon/\sigma^{2}$ and rest length set at the current pair separation.
When run-time bond insertion is required we use ESPResSo~\cite{TODO:espresso} in place of LAMMPS.

\paragraph{Water (Experiment~\ref{subsec:water-phases}).}
Water is modelled with TIP4P/2005~\cite{TODO:tip4p-2005}, with rigid geometry maintained by SHAKE.
We use $|\N_{\text{water}}| = 512$ molecules (1536 atoms) in a cubic box with periodic boundary conditions, a $10$~\AA{} real-space cutoff for the Lennard-Jones contribution, and a particle-particle particle-mesh Ewald treatment of long-range electrostatics with a real-space cutoff of $12$~\AA{}.
The integration timestep is $\delta t = 1$~fs.
The liquid branch is equilibrated in $NPT$ at $P = 1$~atm for $1$~ns and run for $5$~ns of production at each temperature in $\{240, 260, 280, 300, 320, 350, 400, 450, 500\}$~K; the ice branch at $240$~K is initialised from a proton-disordered ice Ih configuration and monitored through a pair correlation function to confirm the crystalline structure is retained.
Velocities are written every $20$~fs.

\paragraph{Molten salt (Experiment~\ref{subsec:molten-salt-ionic-conductivity}).}
Molten NaCl is modelled with the Tosi--Fumi pair potential~\cite{TODO:tosi-fumi} at $216$ ion pairs (432 atoms) in a cubic $NVT$ box set to the experimental density of $1.54$~g\,cm$^{-3}$.
Long-range electrostatics are again treated with PPPM.
The integration timestep is $\delta t = 1$~fs.
At each temperature in $\{1100, 1200, 1400, 1600, 1800\}$~K we equilibrate for $200$~ps and run for $2$~ns of production; velocities are written every $10$~fs.

\paragraph{Reproducibility.}
All input decks, analysis scripts, and random seeds used to generate the figures in Section~\ref{sec:experiments} are available in the repository described in Section~\ref{sec:data-and-availability}; each figure is regenerable from a single command and each run records its seed, source commit, and input hash in a provenance file.

    \section{Experiments}
\label{sec:experiments}

\subsection{Convergence of the Noise Term}
\label{subsec:convergence-of-the-noise-term}

\subsection{Bond Addition}
\label{subsec:bond-addition}

\subsection{Lennard-Jones Phases}
\label{subsec:lennard-jones-phases}
As a first test of the spectral observables introduced in Section~\ref{subsec:interpreting-atomic-correlation-matrices}, we study a single-component Lennard-Jones system swept across its solid, liquid and gas regimes.
Temperatures are reported in Kelvin using argon parameters; the simulation protocol is summarised in Section~\ref{subsec:molecular-dynamics-simulations}.
Figure~\ref{fig:lj-results} compares the mass density (a) with the von Neumann entropy of the atom-wise velocity correlation matrix (b) across the scanned temperature range.
A solid-to-liquid change is visible in both observables near $T \approx 20$--$25$~K, but the entropy shows a noticeably sharper step there than the density.
At $T \gtrsim 85$~K the system leaves the condensed phase, which is reflected as a collapse in both quantities.
\begin{figure}
    \centering
    \includegraphics[width=\linewidth]{figures/lj-experiment.pdf}
    \caption{Lennard-Jones sweep as a function of temperature. (a) Mass density. (b) Von Neumann entropy of the atom-wise velocity correlation matrix $C$ [Equation~\ref{eqn:Cij}]. The solid-to-liquid change near $20$--$25$~K is resolved more sharply by the entropy than by the density, and the loss of the condensed phase above $\sim 85$~K is visible in both.}
    \label{fig:lj-results}
\end{figure}

\subsection{Water Phases}
\label{subsec:water-phases}

\subsection{Molten Salt Ionic Conductivity}
\label{subsec:molten-salt-ionic-conductivity}

    \input{input/conclusion}

\section{Data and Availability}
\label{sec:data-and-availability}


%%%%%%%%%%%%%%%%%%%%%%%%%%%%%%%%%%%%%%%%%%%%%%%%%%%%%%%%%%%%%%%%%%%%%
%% The "Acknowledgement" section can be given in all manuscript
%% classes.  This should be given within the "acknowledgement"
%% environment, which will make the correct section or running title.
%%%%%%%%%%%%%%%%%%%%%%%%%%%%%%%%%%%%%%%%%%%%%%%%%%%%%%%%%%%%%%%%%%%%%
\begin{acknowledgement}

\end{acknowledgement}

%%%%%%%%%%%%%%%%%%%%%%%%%%%%%%%%%%%%%%%%%%%%%%%%%%%%%%%%%%%%%%%%%%%%%
%% The same is true for Supporting Information, which should use the
%% suppinfo environment.
%%%%%%%%%%%%%%%%%%%%%%%%%%%%%%%%%%%%%%%%%%%%%%%%%%%%%%%%%%%%%%%%%%%%%
\begin{suppinfo}

\end{suppinfo}

%%%%%%%%%%%%%%%%%%%%%%%%%%%%%%%%%%%%%%%%%%%%%%%%%%%%%%%%%%%%%%%%%%%%%
%% The appropriate \bibliography command should be placed here.
%% Notice that the class file automatically sets \bibliographystyle
%% and also names the section correctly.
%%%%%%%%%%%%%%%%%%%%%%%%%%%%%%%%%%%%%%%%%%%%%%%%%%%%%%%%%%%%%%%%%%%%%
\newpage
\bibliography{input/bibliography.bib}

\end{document}